\section{Data and evaluation protocol}
\label{sec:protocol}

\subsection{Streams, aggregation, and provenance}

The case study uses public Great Britain electricity-system telemetry from the
National Energy System Operator (NESO). Five demand-side columns are read at
half-hour settlement resolution: national demand (ND), transmission-system
demand (TSD), embedded wind generation, embedded solar generation, and pumped
storage pumping \citep{neso2019historicdemand}. A sixth column is formed from
NESO system-frequency data by taking the maximum absolute deviation from
50~Hz within each settlement period \citep{neso2019systemfrequency}. This
half-hour maximum retains an extreme deviation but not its sub-period timing or
waveform.

The retained 2019 and 2020 inputs are pinned by byte count and SHA-256 checksum.
The report objects identify the generating source files and source-manifest
hashes; the paper tables are generated from those committed reports. The data
are used under the NESO Open Data Licence \citep{nesoOpenLicence}. Required
attribution is: ``Supported by National Energy SO Open Data.''

Rows are ordered by settlement date and period. The same-regime residual
history distinguishes weekday from weekend and settlement period. The loader
handles settlement-period irregularities defensively, while the later alert
budget still groups fixed 48-row blocks; that distinction matters on
daylight-saving transitions.

\subsection{Segments and frozen protocol}

January--April 2019 is designated as the fit segment for the dependence matrix
and AR(1) parameters. It is assumed representative enough for this role; it is
not established to be anomaly-free or ``clean.'' May--June supplies prior
windows before the declared evaluation segment. July--December 2019 is the
development segment, comprising 2,208 scored non-overlapping windows. Choices
of $W=4$, $K=40$, the robust-z marginal, and inclusion of the frequency summary
were informed by 2019 evidence and are not validation outcomes.

The method, sensor set, and configuration were then fixed before the 2020
events were inspected. Rolling references continue across the year boundary,
and all 4,392 scored windows in 2020 form the frozen TRACE-C hold-out. We use
``frozen'' to describe this development chronology, not to assert that the
hold-out is a clean null year. In particular, 2020 contains major weather and
demand-regime changes documented independently
\citep{metoffice2020stormseason,ukgov2020lockdown,ngeso2021endyear}.

\subsection{Event annotations and opportunities}

Annotations are joined only after ranking. The 9 August 2019 power cut and
frequency disturbance is documented by Ofgem \citep{ofgem2020poweroutage}; its
evaluation interval overlaps 5 two-hour windows. Storm Atiyah spans 12 windows,
while the two-day Ciara and Dennis intervals contain 24 windows each, using the
Met Office storm chronology \citep{metoffice2020stormseason}. The first
lockdown evaluation interval is 23 March--5 April 2020, or 168 windows, anchored
to the 23 March announcement \citep{ukgov2020lockdown}. Because event intervals
differ, their best ranks have unequal opportunity for a small value and should
not be treated as like-for-like point-event metrics.

Storm Ellen and Storm Alex were not pre-specified evaluation intervals. Their
names were attached after ranking to the 20 August and 3 October windows using
the external storm record \citep{metoffice2020stormseason,metoffice2020stormalex}.
They are interpretations of unalerted windows, not detector inputs or
discoveries.

\subsection{Post-hoc baselines}

Four baselines share the source rows, six-stream sensor set, $W=4$ windows,
fit cutoff, evaluation dates, score direction, and event intervals. They are a
convolutional autoencoder inspired by standard reconstruction approaches
\citep{sakurada2014autoencoders,vijay2020kerasautoencoder}, PCA reconstruction
with three components \citep{jackson1979pcaresiduals}, a 200-tree Isolation
Forest with fixed seed on window mean/standard-deviation features
\citep{liu2008isolationforest}, and a trailing-context Spectral Residual score
\citep{ren2019spectralresidual}.

Their protocol is explicitly non-identical to TRACE-C. Each baseline uses
global per-stream fit mean/standard-deviation scaling, then directly ranks its
raw anomaly score against a growing strictly-prior history after a 40-window
post-fit warm-up. Selection is an unbudgeted record rule. TRACE-C instead uses
same-regime robust residuals, three trailing 240-window channel ranks, Fisher
aggregation, a growing outer rank, a BH-first selector, and a two-per-block
budget. The baseline suite was constructed after TRACE-C's 2020 results were
known. Thus only TRACE-C was frozen before 2020 inspection, and Table
\ref{tab:baselines} is a post-hoc comparison rather than a second hold-out
experiment.

\subsection{Reproducibility}

The deterministic TypeScript evaluation rebuilds both TRACE-C reports from the
pinned inputs. The baseline score artifact records its fixed random seed,
trainer/source hash, input-series hash, runtime versions, and training history;
the final evaluator verifies those fields before producing the comparison.
Generated \LaTeX{} tables are derived from the JSON evidence rather than edited
manually. These measures support computational reproduction of the reported
arithmetic, while making no claim that different platforms will produce
bitwise-identical external-model artifacts.
